Redes varejistas precisam escolher políticas de reposição para séries de demanda heterogêneas. Esta dissertação investiga se características dessas séries permitem melhorar o compromisso entre custo e atendimento em relação a uma política global. A proposta, denominada \textit{Adaptive Inventory Policy Engine} (AIPE), integra preparação de dados, caracterização da demanda e comparação de políticas em simulação. A avaliação brasileira avalia a coorte de 145 séries em duas origens temporais, com sete políticas incluindo um controle adicional, três sementes nos métodos estocásticos e blocos separados de validação e teste. Os 4.350 ajustes estão completos. A seleção por perfil aumenta o custo médio em 1,93\% e o atendimento em 3,89 pontos percentuais, reduzindo de 119 para 100 os casos abaixo da meta de 0,70 entre 283 casos série--origem com demanda positiva. Os sete casos sem demanda no teste têm atendimento indefinido e são contabilizados separadamente. A escolha global varia entre origens; na segunda, ambos os seletores recomendam a mesma política. A evidência não demonstra economia geral da seleção por perfil, mas explicita o custo do serviço adicional e a sensibilidade temporal das recomendações. A contribuição é a formulação e a avaliação rastreável da seleção contextual no varejo brasileiro. A coorte retrospectiva e as janelas curtas parcialmente sobrepostas limitam a generalização. O componente supervisionado de seleção, \textit{Policy Selection Engine} (PSE), permanece a validar; a avaliação complementar no M5 permanece em andamento.

\vspace{2em}
\textbf{Palavras-chave}: Gestão de estoques; Demanda intermitente; Seleção de políticas; Nível de serviço; Simulação.
